\chapter{RESULTS AND DISCUSSIONS}
\pagebreak


\begin{center}
{\LARGE\textbf{RESULTS AND DISCUSSIONS}}
\end{center}

This chapter presents the results and discussion of the TrapPot implementation and testing. The main objective is to evaluate whether the developed system met its goal of capturing SSH and Telnet interactions, monitoring related network traffic, classifying traffic behavior with the trained Random Forest model, and visualizing the collected data through Kibana.

The analysis treats TrapPot as a complete pipeline rather than a set of separate tools. Cowrie provides honeypot interaction, Zeek provides network monitoring, the Random Forest detector provides traffic classification, and the ELK Stack provides log storage and visualization. The results are therefore discussed according to how well these components worked together during testing.

The chapter is organized into four sections: analysis strategy, analysis results, major findings, and discussion.

\section{Analysis Strategy}
\label{sec:analysis_strategy}

TrapPot was evaluated from both operational and machine learning perspectives. Since the project is a cybersecurity monitoring and detection system, success is not measured only by model accuracy. The full path from honeypot interaction to dashboard visualization must also work correctly.

The analysis was guided by the following questions:

\begin{itemize}
    \item Can Cowrie act as an SSH and Telnet honeypot and log attacker-like interaction?
    \item Can Zeek monitor the traffic generated by the honeypot and produce structured network logs?
    \item Can the AI detector read Zeek's \texttt{conn.log} and classify traffic behavior using the trained Random Forest model?
    \item Can Logstash process the generated logs and send them to Elasticsearch?
    \item Can Kibana display the indexed data through data views and the TrapPot dashboard?
\end{itemize}

The system was tested in a controlled local lab environment using Docker Compose. After the services started, an SSH session was opened to the Cowrie honeypot on port 2222. During the test session, commands such as \texttt{whoami}, \texttt{pwd}, \texttt{ls -la}, \texttt{cat /etc/passwd}, and \texttt{wget} were executed.

These commands were selected because they represent common actions an attacker may take after gaining shell access. For example, \texttt{whoami} checks the active user, \texttt{pwd} checks the current directory, \texttt{ls -la} lists files, \texttt{cat /etc/passwd} attempts to read system account information, and \texttt{wget} simulates a file download attempt.

After the session, each component was checked. Cowrie logs were reviewed to confirm that the login and commands were recorded. Zeek logs were checked to confirm that \texttt{conn.log} and \texttt{ssh.log} were generated. The AI detector output was reviewed to confirm that classification alerts appeared. Elasticsearch indices were checked to confirm that logs were indexed, and Kibana was opened to confirm that the data views and dashboard were available.

\section{Analysis Result}
\label{sec:analysis_result}

The test results showed that the main parts of TrapPot worked as a connected Docker-based system. Docker Compose started the required services, including Cowrie, Zeek, the AI detector, Elasticsearch, Logstash, Kibana, and the setup services. This confirmed that the project can be deployed as a portable multi-container environment.

Cowrie successfully accepted an SSH connection on port 2222. After login, the honeypot provided a fake Linux shell environment and accepted the commands entered during the session. Cowrie recorded the login attempt, successful authentication, executed commands, and file download attempt. This result shows that the honeypot layer achieved its purpose of attracting and capturing attacker-like behavior without exposing a real production system.

Zeek monitored the traffic generated by the Cowrie session. The connection records in \texttt{conn.log} showed SSH traffic between the host machine and the Cowrie container. Zeek also captured traffic generated by the simulated \texttt{wget} command. This confirms that the network monitoring layer converted traffic activity into structured log records.

The AI detector read Zeek's \texttt{conn.log} and produced live classification alerts. During testing, the detector produced output such as \texttt{Normal detected}. This validates that the detector could extract the relevant network features, load the trained Random Forest model, and produce a runtime prediction. The detector also wrote detection records to \texttt{detections.json}, allowing the detection output to be stored and later indexed.

Logstash and Elasticsearch were also validated. Logstash connected to Elasticsearch and parsed existing TrapPot logs. Elasticsearch contained the expected TrapPot indices for Cowrie, Zeek, and detector output. This confirms that the log processing and indexing layer worked correctly.

Kibana was available at \texttt{localhost:5601}. Data views were available for Cowrie logs, Zeek connection logs, Zeek SSH logs, and AI detector results. The \texttt{TrapPot Overview} dashboard was also available. This validated that the visualization layer was configured correctly and could be used to inspect collected logs through a browser interface.

Table \ref{tab:component_verification} summarizes the main component verification results.

\begin{table}[htbp]
    \centering
    \caption{TrapPot Component Verification Results}
    \label{tab:component_verification}
    \begin{tabular}{|p{4cm}|p{7cm}|p{2cm}|}
        \hline
        \textbf{Component} & \textbf{Verification Method} & \textbf{Result} \\
        \hline
        Cowrie & SSH login was performed on port 2222 and commands were executed in the fake shell. & Passed \\
        \hline
        Zeek & \texttt{conn.log} and \texttt{ssh.log} were generated after the SSH test session. & Passed \\
        \hline
        AI Detector & The detector read Zeek's \texttt{conn.log} and produced live classification alerts. & Passed \\
        \hline
        Logstash & Logstash processed TrapPot logs and connected to Elasticsearch. & Passed \\
        \hline
        Elasticsearch & TrapPot-related Cowrie, Zeek, and detector indices were created. & Passed \\
        \hline
        Kibana & Data views and the \texttt{TrapPot Overview} dashboard were available in the browser. & Passed \\
        \hline
    \end{tabular}
\end{table}

The Random Forest model was tested on the prepared IoT-23 dataset \cite{iot23dataset}. The final trained model achieved approximately 99.59\% accuracy across the five selected classes. This result indicates that the model learned the prepared dataset effectively, although live deployment performance should be validated further with larger volumes of real honeypot traffic.

Table \ref{tab:model_accuracy} presents the model accuracy result.

\begin{table}[htbp]
    \centering
    \caption{Random Forest Model Accuracy Result}
    \label{tab:model_accuracy}
    \begin{tabular}{|p{6cm}|p{4cm}|}
        \hline
        \textbf{Model} & \textbf{Accuracy} \\
        \hline
        Random Forest Classifier & 99.59\% \\
        \hline
    \end{tabular}
\end{table}

\section{Major Findings}
\label{sec:major_findings}

The first major finding is that TrapPot integrated honeypot interaction, network monitoring, machine learning detection, and dashboard visualization into one working environment. The system did not only collect terminal logs. It also transformed network activity into structured logs and used those logs for live classification.

The second finding is that Zeek is important in linking the honeypot to the AI detector. Cowrie records attacker behavior at the interaction level, such as login attempts and commands. Zeek records network activity, including ports, protocols, services, duration, byte counts, packet counts, and connection state. This makes the detector more useful because it classifies traffic using structured network features rather than raw command logs alone.

The third finding is that the AI detector was successfully connected to the live system. The detector loaded \texttt{trappot\_model.pkl} and \texttt{label\_encoder.pkl}, monitored Zeek's \texttt{conn.log}, extracted the twelve required features, and produced live classification alerts. This confirms that the machine learning component was integrated with the implemented pipeline and was not only an offline model.

The fourth finding is that Docker Compose improved deployment and portability. A single Docker Compose workflow can start the full system instead of requiring each component to be installed and configured manually.

The fifth finding is that Kibana made the system easier to inspect and present. Terminal logs are useful for debugging, but Kibana provides a clearer interface for searching, filtering, and viewing activity through dashboards.

\section{Discussion}
\label{sec:discussion}

The results show that TrapPot achieved its main implementation goals. The system recorded SSH honeypot activity, generated network logs with Zeek, produced AI classification alerts, stored logs in Elasticsearch, and visualized the data in Kibana. This shows that the proposed architecture can function as a honeypot-based monitoring and detection system.

A key strength of the system is that it combines two levels of observation. Cowrie records the attacker's interaction with the fake system, while Zeek records the related network traffic. This combination provides both behavioral and network-level views of the same activity. The AI detector then uses Zeek's network features to classify traffic behavior.

There are still limitations. First, the system was tested in a controlled local environment rather than on a public network. This was necessary for safety and ethical reasons, but it means the system was not exposed to large volumes of real attacker traffic. Future testing should be conducted in an approved isolated lab network to capture more realistic traffic.

Second, the live test activity was limited. The SSH session was manually generated and small in scale, so it may not represent the full range of attack patterns in the IoT-23 training data. More varied attack simulations are needed to evaluate the model under different live conditions.

Third, the Random Forest model achieved high accuracy on the prepared IoT-23 dataset, but dataset accuracy does not guarantee the same performance in a live deployment. Real environments may contain traffic patterns that differ from the training data. Further validation using live honeypot traffic and additional attack scenarios is therefore required.

Fourth, Elasticsearch, Logstash, and Kibana consume more resources than the core honeypot and detector components. They provide strong storage and visualization capabilities, but long-term monitoring or higher traffic volume may require memory and performance tuning.

Overall, the results show that TrapPot is a successful proof-of-concept system. It demonstrates how Cowrie, Zeek, Random Forest classification, and the ELK Stack can be integrated in a Docker-based environment for IoT honeypot monitoring and analysis. The system provides a strong foundation for future improvements, including broader attack testing, expanded validation, and more detailed dashboard analytics.
